% ============================================================
%  Bonus, Segment 3A — Mastering the Fretboard / 指板精通
%  原书页码 74-82
% ============================================================
\chapter{Mastering the Fretboard / 指板精通}
\label{ch:fretboard}

\eng{There are many methods out there to help you become familiar with note names and their locations on the neck: arpeggios, scales, movable chord shapes, chord shapes starting from different strings, the list goes on. I encourage you to try out different methods to find what works best for you. The following are methods I found rather useful for finding my way around the guitar neck more easily, more logically, and helped me on my way to mastering the fretboard.}

\chn{有很多方法可以帮助你熟悉音名及其在指板上的位置：琶音、音阶、可移动和弦形状、从不同弦开始的和弦形状，等等。我鼓励你尝试不同的方法，找到最适合你的。下面是我认为非常有用的一些方法，能够帮助你更轻松、更有逻辑地掌握指板，走向指板精通。}

\paragraph{Finding Notes Quickly — The Octave Method}
\eng{This method is a useful way to quickly locate notes all over the fretboard. It proves valuable for identifying root notes in chord shapes, arpeggios, scales, and more. I call it the "octave method," although it might have another name. Initially, focus on learning the positions of notes on the Low E and A strings (in standard tuning). While this may seem daunting if you are new to it, with practice, it becomes less challenging than you might expect. I have listed the notes starting from the open E and A strings up to their respective octaves on the twelfth fret (beyond the twelfth fret, you'll find the same notes but an octave higher).}

\chn{这种方法是快速定位指板上音符的实用方式。它在识别和弦形状、琶音、音阶等的根音时非常有价值。我称它为“八度法”，虽然它可能还有其他名称。最初，重点学习标准调弦中低音 E 和 A 弦上的音符位置。对于新手来说这可能看起来有些困难，但通过练习，它会比你想象的容易。我列出了从空弦 E 和 A 到十二品对应八度的音符（超过十二品后，你会找到相同的音，但高一个八度）。}

\eng{To illustrate how to apply this concept, I have provided a diagram. I have marked the note G on the low E and A strings with blue circles. Additionally, I have indicated the octaves of these notes (the same note but higher in pitch) with black circles. This can help you visualize how notes and their octaves are located on the guitar neck.}

\chn{为了说明如何应用这个概念，我提供了一个示意图。我用蓝色圆圈标出了低音 E 和 A 弦上的 G 音，并用黑色圆圈标出了这些音的八度（相同音但更高）。这有助于你可视化音符及其八度在指板上的位置。}

\eng{Notice how they have a particular spacing between each octave. From the G note (third fret) on the low E string, we jump a string (the A string) and find the octave G three frets higher (fifth fret on the D string). The same applies for the note G played on the A string. You may be familiar with these octave shapes if you have ever learned a pop punk or punk song. However, the spacing does change slightly when we start from the G note on the D string, and the same for the G note starting from the G string, the spacing is now four frets instead of three. The orange notes in the example below highlight this. I would not worry about getting this to stick in your mind because when played correctly you will hear them as octaves. If you play three fret spaces instead of four for octaves starting on the D and G strings you will get this major seven interval that clashes in sound, so you will instantly realize you need to slide your higher note up one fret to get the octave sound.}

\chn{注意每个八度之间有特定的间距。从低音 E 弦上的 G 音（第三品）开始，我们跳一根弦（A 弦），在 D 弦第五品找到高八度的 G。A 弦上的 G 也有同样的关系。如果你学过流行朋克或朋克歌曲，你可能已经熟悉这些八度形状。然而，当我们从 D 弦上的 G 开始时，间距会略有变化；从 G 弦开始时，间距变为四品而不是三品。下面示例中的橙色音符强调了这一点。不必过分担心记忆这些形状，因为正确演奏时你会听到它们是八度。如果你在 D 和 G 弦上以三品而不是四品来演奏八度，就会出现不协调的增七音程，你会立刻意识到需要把高音上移一个品位才能得到八度音。}

\eng{Another handy trick is to notice that the low E and high E strings follow the same note order. Meaning you can mirror any note you play on either string to find the same note. I have indicated this with the red notes on the chart below.}

\chn{另一个实用技巧是注意低音 E 弦和高音 E 弦具有相同的音序。这意味着你可以将任意一个 E 弦上的音在另一根 E 弦上镜像找到相同的音。我在下面的图表中用红色音符标出了这一点。}

\eng{This system, I believe, is suited to players who have yet to have a built-up knowledge of note locations on the fretboard, so by following this method, you will slowly but surely start building your knowledge of note names and where they are located across the fretboard.}

\chn{我认为这个系统适合那些还没有积累指板音符位置知识的玩家。通过这种方法，你会逐步建立起音名及其在指板上位置的认知。}

\paragraph{Visualizing the Neck in Shapes - The Arpeggio Method}
\eng{By learning this method, you will be able to conjure up something both quickly and easily anywhere on the neck. Even something to play over fast changing or complex chord changes, and if time permitted you can really write something you will be proud of. Not only will it help you achieve the former, but it will also provide another useful approach to writing riffs, licks, melody lines, and so on. How and why this method will help you will become clearer as you work through this section.}

\chn{通过学习这种方法，你将能够在指板任何位置快速且轻松地演奏东西。即使是在快速变化或复杂和弦变化之上也能奏出合适的内容，如果时间允许，你甚至可以真正写出令自己满意的乐句。它不仅帮助你实现上述目标，还为写作 riff、lick、旋律线等提供了另一种有用的方法。随着你阅读本节，这种方法如何帮助你以及为什么如此有用会变得更加清晰。}

\eng{In this section, four types of extended arpeggios will be presented: major seven, minor seven, dominant, and minor seven flat five. Compared to arpeggio shapes you may have already learned; I will be taking a slightly different approach by teaching you arpeggios shapes that combine to form a pattern across the fretboard that you can easily learn. Also, I have included examples of how you could apply them.}

\chn{本节将介绍四种扩展琶音：大七、小七、属七和半减七。与您可能已经学过的琶音形状相比，我将采用略有不同的方法，教你一些可以组合成指板上可学习模式的琶音形状。同时，我还会提供如何应用它们的示例。}

\eng{Extended seven arpeggios give us an extra note to work with when compared with basic chord arpeggios (refer to the theory section for more detailed information) and will 'fit' over any iteration of a particular unaltered chord type. For example, a C major 7 arpeggio will work over a C major chord, C sus 2 or sus 4 chord, a C major 6 chord, C major 7 chord, C major 9 chord, and so on. If the chord or implied harmony is not altered, you should be safe to use it. A Cmaj7b5 chord is an example of an altered chord, the flattened 5th changes the quality of the chord, meaning it would NOT be harmonically sound if you were to play a Cmaj7 arpeggio along with it. Chords that contain notes such as b9, \#9, b5 \#5 are referred to as being altered chords, there is a lot to unpack there, but just think of it as making a chord sound more 'out there' in terms of sound.}

\chn{与基本和弦琶音相比，扩展七琶音提供了额外的音符（详见乐理章节），并且可以适用于任何未改变性质的相同和弦类型。例如，C 大七琶音可用于 C 大和弦、Csus2、Csus4、C6、Cmaj7、Cmaj9 等。如果和弦或隐含和声没有被改变，你使用它通常是安全的。Cmaj7b5 是一个被改变的和弦示例，降五度改变了和弦的性质，这意味着如果你在它上面演奏 Cmaj7 琶音，就不会和谐。包含 b9、\#9、b5、\#5 等音的和弦被称为改变和弦，这里面内容很多，但你可以把它理解为让和弦听起来更“异样”的一种方式。}

\paragraph{Major Seven Arpeggios Across the Fretboard}
\eng{Below you will find major seven arpeggios spanning one octave, meaning that when you get to the twelfth fret on your guitar, you should repeat the same patterns again. I have colour coded each shape to show how they are connected to each other and so you can see each individual arpeggio (chord) shape. Take your time working through each shape and find the most comfortable fingering, next, experiment with stringing them together into a chain, like in the transcribed version below, to help ingrain them into your memory and muscle memory.}

\chn{下面你会看到跨一个八度的大七琶音，也就是当你到达第十二品时，应重复相同模式。我用颜色标识了每个形状，展示它们之间的连接，以便你看清每个单独的琶音（和弦）形状。慢慢练习每个形状，找到最舒适的指法，然后尝试将它们串联成一个链条，就像下面的转录版本那样，帮助你将它们记住并形成肌肉记忆。}

\eng{Here is how they look transcribed for better clarity.}

\chn{下面是它们的转录版本，便于更清晰理解。}

\eng{Each individual chord shape used for the arpeggios are listed below. You can use any of these chord shapes interchangeably for major seven chords. Meaning, you will have to adjust the shapes when starting from a different root note, but the shapes will remain the same.}

\chn{下面列出了每个琶音所使用的和弦形状。你可以将这些形状互换用于大七和弦。这意味着当你从不同根音开始时需要调整位置，但形状本身保持不变。}

\eng{The following are a couple of examples of ideas you could potentially write. You could extend the examples here, but you will soon run into problems of things sounding too similar, thus these arpeggios are best suited to being combined with other arpeggios, scales and so on. At the end of this section, I will show examples of how you can combine a chain of arpeggios together for some musical phrases.}

\chn{下面是几个你可能创作的示例。你可以延展这些示例，但很快你会发现听起来太过相似，因此这些琶音最适合与其他琶音、音阶等结合使用。本节末我会展示如何将一串琶音组合成音乐短句的示例。}

\paragraph{Minor Seven Arpeggios Across the Fretboard}
\eng{Moving on to minor seven arpeggios, again they span one octave, meaning when you get to the twelfth fret, repeat the same patterns again. I have colour coded each shape to show how they are connected and so you can see each individual arpeggio (chord) shape. Take your time to work through each shape and find the most comfortable fingering for you, next, experiment with stringing them together into a chain like in the transcribed version below to help ingrain them into your memory and muscle memory.}

\chn{接下来是小七琶音，同样跨一个八度，也就是到达第十二品时重复相同模式。我用颜色标识了每个形状，展示它们之间的连接，以便你看清每个单独的琶音（和弦）形状。慢慢练习每个形状，找到最舒适的指法，然后尝试将它们串联成一个链条，就像下面的转录版本那样，帮助你将它们记住并形成肌肉记忆。}

\paragraph{Dominant Arpeggios Across the Fretboard}
\eng{Next, we have dominant arpeggios spanning one octave. Again, I have colour coded each shape to show how they are connected so you can see each individual arpeggio (chord) shape. Take your time to work through each shape and find the most comfortable fingering for you.}

\chn{接下来是属七琶音，同样跨一个八度。我同样用颜色标识了每个形状，展示它们之间的连接，以便你看清每个单独的琶音（和弦）形状。慢慢练习每个形状，找到最舒适的指法。}

\eng{Again, here are a few examples of short ideas you could potentially write. They sound a little odd in isolation, but they do work as individual licks or riff ideas. They sound better with the relevant dominant chord behind them (C7 for these examples), and the perhaps the following examples feel limited because they would work better in conjunction with a changing chord progression, meaning the arpeggios would change too. However, this is math rock so do whatever suits best!}

\chn{再次给出几个你可能创作的简短示例。它们单独听起来可能有点奇怪，但作为独立的乐句或riff想法是可行的。它们与相应的属七和弦（例如本例中的 C7）搭配时会更好。或许下面的示例听起来有限，因为它们更适合与变化的和弦进行配合使用，也就是说琶音也会随之变化。不过这毕竟是数学摇滚，按你觉得合适的方式去做吧！}

\paragraph{Minor Seven Flat Five (Half Diminished) Arpeggios Across the Fretboard}
\eng{Lastly, let us conclude with minor seven flat five arpeggios spanning one octave. Again, I have colour coded each shape to show how they are connected so you can see each individual arpeggio (chord) shape. However, I have changed the root note from C to D to better fit them in the examples. Take your time to work through each shape and find the most comfortable fingering for you. From there, experiment with stringing them together into chunks. The transcribed version below will aid in ingraining them into your memory and muscle memory.}

\chn{最后，我们以跨一个八度的半减七琶音结束。我同样用颜色标识了每个形状，展示它们之间的连接，以便你看清每个单独的琶音（和弦）形状。不过我把根音从 C 改为 D，以便示例更契合。慢慢练习每个形状，找到最舒适的指法。然后尝试将它们串成片段。下面的转录版本将帮助你将它们记住并形成肌肉记忆。}

\paragraph{Arpeggios - Examples Riffs \& Licks}
\eng{Before moving forward, it is important that you are already familiar with the extended arpeggios we have just explored. In the upcoming section, the goal is to provide you with a few examples of how to combine different types of arpeggios to create riffs and licks. These can be played alongside a chord progression or used independently. These examples are intentionally designed to be relatively simple as this will make it easier for you to learn and grasp the concepts. As you work through them, try to identify which arpeggio shapes are being used.}

\chn{在继续之前，重要的是你已经熟悉我们刚才探讨的扩展琶音。下一个部分的目标是为你提供一些示例，展示如何组合不同类型的琶音来创造 riff 和 lick。这些示例可以与和弦进行一起演奏，也可以独立使用。这些示例特意设计得相对简单，以便你更容易学习和掌握概念。练习时，尝试识别正在使用哪些琶音形状。}

\eng{The following is an example of how you could write some ascending style licks. To make things
easier notice the triplet groupings (grouping of three notes). Seeing it as groupings of three
notes it will make it much easier on the fingers and your mind. Try them backwards too, for a
descending feel.}

\chn{下面是一个你可以写成上行风格乐句的示例。为了更容易理解，请注意三连音分组（每组三个音）。将其视为三音一组会让你的手指和思维更容易驾驭。也尝试反过来演奏，获得下行感觉。}

\eng{Example 4G is a rather crude one, but I designed it to get you bouncing around the arpeggio shapes in two note groupings. Can you work out which arpeggio shapes the note groupings are derived from? Noticing them will help cement ideas in your memory better.}

\chn{示例 4G 有些粗糙，但我设计它是为了让你在两音分组中在琶音形状间跳跃。你能找出这些音组来源于哪些琶音形状吗？注意它们会帮助你更好地将概念记住。}

\eng{Example 4H takes the same idea but hashes it out into a passage. Sounds much more musical, right?}

\chn{示例 4H 采用相同的想法，但将其扩展成一段乐句。听起来更富音乐性了，对吧？}

\eng{Example 4I is another more musical idea. It is based around a IV, V, vi progression. You will hear the band Delta Sleep use this kind of chord progression quite regularly and to great effect in their music.}

\chn{示例 4I 是另一个更具音乐性的想法。它基于 IV、V、vi 的进行。你会发现 Delta Sleep 乐队经常使用这种和弦进行，并且效果极佳。}

\eng{Example 4J is another mathy riff sounding riff based round a I-IV progression. This time, I wanted to have quicker changes and some odd timing to make things more musical and an example of what you are likely to see in math rock styles of guitar. If you count the first bar as a grouping of three and four beats, and the second bar as three groupings of three it will make things easier.}

\chn{示例 4J 是另一个具有数学摇滚风格的 riff，基于 I-IV 进行。这一次，我希望加入更快的变化和一些奇拍，以使乐句更具音乐性，并展示你在数学摇滚吉他风格中可能会看到的内容。如果你把第一小节看作 3+4 拍分组，第二小节看作三个 3 拍分组，会更容易理解。}

\eng{Lastly, I could not help but include a couple of tapping examples. This is math rock after all!}

\chn{最后，我忍不住也要包含几个点弦示例。毕竟这是数学摇滚嘛！}

\eng{As time passes, and the more time you dedicate to these methods, you will better internalize where notes are on the fretboard and formulate your ideas to greater effect. This concluded this segment. Next, we will look at functional strumming patterns.}

\chn{随着时间推移，你对这些方法投入的时间越多，就会越能内化指板上音符的位置，并更有效地形成创意。至此，本节结束。接下来，我们将看看功能性扫弦模式。}
